% Teaching construction: pi_1 = pi_2 = 1/2; mu_1 = (-1,0), mu_2 = (1,0).
% Left: Sigma_1 = Sigma_2 = I. Right: Sigma_1 = I, Sigma_2 = diag(1,4).
% Contours are Delta_c = 1, not equal density values across classes.
% Right boundary: g_2-g_1 = 2*x_1 + 3*x_2^2/8 - ln(2) = 0.
\begin{tikzpicture}[
  x=0.072\linewidth,y=0.072\linewidth,
  font=\footnotesize,line cap=round
]
  \foreach \shift in {0,0.5} {
    \begin{scope}[xshift=\shift\linewidth]
      \draw[book arrow] (-2.5,-2.5) -- (2.7,-2.5);
      \draw[book arrow] (-2.5,-2.5) -- (-2.5,2.65);
      \node at (2.6,-3.05) {$x_1$};
      \node at (-2.8,2.85) {$x_2$};
      \foreach \tick in {-2,0,2} {
        \draw[FigureInk] (\tick,-2.5) -- (\tick,-2.58);
        \node[anchor=north] at (\tick,-2.6) {$\tick$};
        \draw[FigureInk] (-2.5,\tick) -- (-2.58,\tick);
        \node[anchor=east] at (-2.6,\tick) {$\tick$};
      }
      \draw[FigureBlue,line width=1.2pt] (-1,0) circle[radius=1];
      \fill[FigureBlue] (-1,0) circle[radius=2.2pt];
      \path[fill=FigureYellow] (1,0) ++(-2.2pt,-2.2pt)
        rectangle ++(4.4pt,4.4pt);
      \node[text=FigureBlue] at (-1,-0.5) {$\bm{\mu}_1$};
      \node[text=FigureYellow] at (1,-0.5) {$\bm{\mu}_2$};
    \end{scope}
  }

  \draw[FigureYellow,dashed,line width=1.2pt] (1,0) circle[radius=1];
  \draw[FigureRed,line width=1.25pt] (0,-2.5) -- (0,2.5);
  \node[font=\figurefont\footnotesize] at (0,3.65) {共享协方差};
  \node at (0,3.1) {$\bm{\Sigma}_1=\bm{\Sigma}_2=\bm{I}_2$};
  \node[text=FigureInk,inner sep=1pt] at (0,-3.7)
    {\textbf{LDA：}$x_1=0$（线性）};

  \begin{scope}[xshift=0.5\linewidth]
    \draw[FigureYellow,dashed,line width=1.2pt]
      (1,0) ellipse[x radius=1,y radius=2];
    \draw[FigureRed,line width=1.25pt]
      plot[domain=-2.5:2.5,samples=101,variable=\t]
        ({ln(2)/2-3*\t*\t/16},{\t});
    \node[font=\figurefont\footnotesize] at (0,3.65) {类特有协方差};
    \node at (0,3.1) {$\bm{\Sigma}_2=\operatorname{diag}(1,4)$};
    \node[text=FigureInk,inner sep=1pt] at (0,-3.7)
      {\textbf{QDA：}$x_1=\frac12\log2-\frac{3}{16}x_2^2$（二次）};
  \end{scope}
\end{tikzpicture}
